\section{应用层}

\subsection{域名系统 DNS}

\begin{mydef}{域名系统（DNS）}{dns-def}
DNS（Domain Name System）是互联网的命名系统，将便于记忆的域名转换为机器使用的 IP 地址。DNS 采用\textbf{客户/服务器（C/S）}模型，运行在 UDP 的 53 号端口上（数据量小、要求快速响应）。当数据超过 512 字节或区域传送时使用 TCP。
\end{mydef}

\subsubsection{域名空间}

\begin{conclusion}{域名的层次结构}{dns-hierarchy}
域名采用层次树状结构，从右向左层次递增，各级之间用「.」分隔：
\[
\texttt{www.example.com.}
\]
最右边有隐含的根域「\texttt{.}」：
\begin{itemize}
    \item \textbf{根域}：用「\texttt{.}」表示，由 13 组根域名服务器管理。
    \item \textbf{顶级域（TLD）}：\texttt{.com}（商业）、\texttt{.org}（非营利组织）、\texttt{.net}（网络服务）、\texttt{.cn}（中国）、\texttt{.edu}（教育）等。
    \item \textbf{二级域（SLD）}：\texttt{example}（由组织注册）。
    \item \textbf{三级域/子域}：\texttt{www}（具体主机名）。
\end{itemize}
域名不区分大小写，每级不超过 63 字符，总长不超过 255 字符。
\end{conclusion}

\begin{attention}
\textbf{408 常考辨析：} 域 vs 区（Zone）。区是域名空间中一个服务器的管辖范围；一个域可能包含多个区，但一个区只属于一个域。区的概念用于 DNS 服务器的实际管理和区域传送。
\end{attention}

\subsubsection{域名服务器}

\begin{conclusion}{DNS 服务器四层结构}{dns-servers}
\begin{tablebox}{DNS 服务器的层次结构}
\centering
\begin{tabularx}{\linewidth}{|p{0.27\linewidth}|X|p{0.27\linewidth}|}
\hline
\textbf{类型} & \textbf{作用} & \textbf{管辖范围} \\
\hline
根域名服务器 & 返回顶级域服务器的 IP & 全球 13 组（a.root-servers.net $\sim$ m）\\
顶级域服务器 & 返回权限域名服务器的 IP & .com, .org, .cn 等 \\
权限域名服务器 & 返回域名—IP 的最终映射 & 特定组织/机构 \\
本地域名服务器 & 代理主机进行查询 & ISP/局域网（非层次结构）\\
\hline
\end{tabularx}
\end{tablebox}
\textbf{补充：} 本地域名服务器不属于 DNS 层次结构，但每个 ISP/大学/公司都有本地域名服务器，距离主机很近（通常在同一局域网）。
\end{conclusion}

\subsubsection{DNS 解析过程 — 递归查询 vs 迭代查询}

\begin{conclusion}{DNS 查询的两种模式}{dns-resolution}
\textbf{递归查询（Recursive Query）：} 被请求的 DNS 服务器必须向请求者返回最终结果（或失败），而不能只给出下一台服务器的转介信息；若它没有答案，就代表请求者继续向下游服务器查询。主机的存根解析器通常向本地 DNS 发递归请求。若后续各级也都采用递归，则查询逐级向下传递、结果逐级返回。

完全递归的消息方向为：主机请求本地 DNS，本地 DNS 请求根 DNS，根 DNS 请求顶级 DNS，顶级 DNS 请求权威 DNS；查询结果再按权威 DNS、顶级 DNS、根 DNS、本地 DNS、主机的方向逐级返回。

\textbf{迭代查询（Iterative Query）：} 当被查询的服务器没有所需记录时，\textbf{它返回下一步应查询的服务器的 IP 地址}，由请求者自己继续查询。本地域名服务器向根、顶级、权限服务器的查询通常为迭代。

典型的迭代流程为：主机递归请求本地 DNS；本地 DNS 查询根 DNS，得到顶级 DNS 地址；再查询顶级 DNS，得到权威 DNS 地址；最后查询权威 DNS，得到最终记录并返回主机。
\end{conclusion}

\begin{attention}
\textbf{408 重难点：递归 vs 迭代的辨析。}
\begin{itemize}
    \item \textbf{常用模式：} 主机的存根解析器 $\to$ 本地域名服务器通常采用\underline{递归查询}，主机不承担逐级查询负担。
    \item \textbf{可配置部分：} 本地域名服务器 $\to$ 其他服务器，可以是递归也可以是迭代。\textbf{实际互联网通常采用：} 主机递归 $\to$ 本地 DNS，本地 DNS \underline{迭代} $\to$ 根/顶级/权限。
    \item 关键字：\textbf{递归 = 全套代劳}（给我结果，别让我再去问）；\textbf{迭代 = 指路}（我不知道，但你去问 XXX）。
    \item 递归查询中每一层都向上一层递归请求，导致根服务器压力巨大，实际中很少使用从本地到根的递归。
\end{itemize}
\end{attention}

\subsubsection{DNS 缓存}

\begin{conclusion}{DNS 缓存机制}{dns-cache}
为提高效率、减轻根服务器压力，DNS 广泛使用缓存：
\begin{itemize}
    \item 本地域名服务器获得一条映射后，将其\textbf{缓存}一段时间（由 TTL 字段决定），后续相同查询直接返回。
    \item 主机本身也可能在浏览器或操作系统的存根解析器中缓存 DNS 记录；\texttt{hosts} 文件是本地静态名称映射来源，不属于 DNS 缓存。
    \item \textbf{TTL（生存时间）}：由权限域名服务器设置，默认 TTL 到期后缓存条目被清除，需要重新查询。
    \item 缓存可能导致不一致：当 IP 地址更改时，在 TTL 过期前缓存中的旧记录可能仍然生效。
\end{itemize}
\end{conclusion}

\subsection{文件传输协议 FTP}

\begin{mydef}{FTP 协议}{ftp-def}
FTP（File Transfer Protocol）用于在网络上进行文件传输，采用 C/S 模型，提供交互式访问。FTP 使用 TCP 协议保证可靠传输。
\end{mydef}

\subsubsection{控制连接与数据连接}

\begin{conclusion}{FTP 的双连接机制}{ftp-double-connection}
FTP 的独到之处在于它使用\textbf{两个并行的 TCP 连接}：
\begin{tablebox}{FTP 控制连接 vs 数据连接}
\centering
\begin{tabularx}{\linewidth}{|p{3.4cm}|X|X|}
\hline
 & \textbf{控制连接} & \textbf{数据连接} \\
\hline
端口号 & 服务器端口 21 & 服务器端口 20 \\
持续时间 & 整个会话期间保持 & 每次文件传输时建立 \\
传输内容 & FTP 命令与应答 & 文件数据/目录列表 \\
协议 & 客户端发命令，服务器应答 & 直接传输数据 \\
\hline
\end{tabularx}
\end{tablebox}

\textbf{控制连接}：客户端通过端口 21 向服务器发起连接，发送 FTP 命令（USER, PASS, LIST, RETR, STOR 等）及接收服务器应答。

\textbf{数据连接}：每当需要传输文件内容或目录列表时，动态建立一条新的数据连接。数据传输完成后，数据连接关闭，控制连接继续保持。

这种「带外」（Out-of-Band）控制方式是 FTP 区别于其他协议（如 HTTP 把命令和数据混在同一条连接中）的核心特征。
\end{conclusion}

\subsubsection{主动模式 vs 被动模式}

\begin{conclusion}{FTP 两种数据连接模式}{ftp-mode}
根据数据连接建立的发起方不同，FTP 分为两种工作模式：

\textbf{主动模式（Active Mode / PORT）：}
\begin{itemize}
    \item 客户端发送 PORT 命令，告知服务器自己在哪个端口（$N > 1023$）上监听。
    \item \textbf{服务器主动}从端口 20 向客户端的 $N$ 端口发起 TCP 连接。
    \item \textbf{问题：} 客户端通常在防火墙/NAT 后面，服务器反向主动连接常被防火墙拦截。
\end{itemize}

\textbf{被动模式（Passive Mode / PASV）：}
\begin{itemize}
    \item 客户端发送 PASV 命令，服务器响应一个随机端口号 $M > 1023$。
    \item \textbf{客户端主动}向服务器的 $M$ 端口发起数据连接。
    \item \textbf{优点：} 防火墙友好——所有连接都由客户端发起。现代 FTP 客户端默认使用被动模式。
\end{itemize}

\begin{attention}
\textbf{408 常考：主动/被动模式的区分。}
\begin{itemize}
    \item \textbf{主动模式}：服务器从 20 端口 $\to$ 客户端端口（服务器主动发起数据连接）——属于从外到内的连接。
    \item \textbf{被动模式}：客户端 $\to$ 服务器随机端口（客户端主动发起数据连接）——所有连接方向一致。
    \item 区分关键：看数据连接由谁发起。无论哪种模式，控制连接始终由客户端发起（到 21 端口）。
\end{itemize}
\end{attention}
\end{conclusion}

\subsection{电子邮件}

\begin{mydef}{电子邮件系统体系结构}{email-arch}
电子邮件系统由三部分组成：
\begin{enumerate}
    \item \textbf{用户代理（UA，User Agent）}：用户与邮件系统的接口（Outlook, Thunderbird, Web Mail）。
    \item \textbf{邮件服务器（Mail Server）}：存储和转发邮件。
    \item \textbf{邮件协议}：SMTP（发送）、POP3/IMAP（接收）。
\end{enumerate}
邮件的整个生命周期为：发送方 UA 使用 SMTP 向发送方邮件服务器提交邮件；发送方邮件服务器使用 SMTP 转交给接收方邮件服务器；接收方 UA 再通过 POP3 或 IMAP 读取邮件。
\end{mydef}

\subsubsection{SMTP — 简单邮件传送协议}

\begin{conclusion}{SMTP 协议}{smtp}
SMTP（Simple Mail Transfer Protocol）是\textbf{发送}邮件的协议。服务器间转发通常使用 TCP 端口 25；客户端邮件提交通常使用端口 587，并可通过 STARTTLS 加密（端口号本身不等于已经加密）。

\begin{itemize}
    \item \textbf{工作过程}：使用命令/应答方式。客户端发送命令（HELO, MAIL FROM, RCPT TO, DATA, QUIT），服务器应答（$220, 250, 354, 221$ 等）。
    \item \textbf{报文格式}：信封（envelope）+ 内容。信封 = MAIL FROM + RCPT TO；内容 = 首部 + 空行 + 主体。
    \item \textbf{首部字段}：\texttt{From:}（邮件是谁写的），\texttt{To:}（发给谁），\texttt{Subject:}（主题）等。注意：这些是邮件首部的文本内容，不是 SMTP 命令。
    \item 传统 SMTP 的邮件数据按 7 位 ASCII 设计，二进制附件和非 ASCII 内容通常借助 MIME 的内容类型与传输编码表示；扩展 SMTP 还定义了 8BITMIME 等能力。
    \item SMTP 使用\textbf{持久连接}：同一连接中可以发送多封邮件。
\end{itemize}

\textbf{典型 SMTP 会话示例：}
\begin{lstlisting}
S: 220 smtp.example.com
C: HELO client.com
S: 250 Hello
C: MAIL FROM: <alice@client.com>
S: 250 OK
C: RCPT TO: <bob@server.com>
S: 250 OK
C: DATA
S: 354 Start mail input
C: From: alice@client.com
C: To: bob@server.com
C: Subject: Hello
C:
C: This is the message body.
C: .
S: 250 OK
C: QUIT
S: 221 Bye
\end{lstlisting}
\end{conclusion}

\subsubsection{POP3 与 IMAP — 邮件读取协议}

\begin{conclusion}{POP3 vs IMAP}{pop3-imap}
\begin{tablebox}{POP3 与 IMAP 对比}
\centering
\begin{tabularx}{\linewidth}{|p{3.4cm}|X|X|}
\hline
 & \textbf{POP3} & \textbf{IMAP} \\
\hline
TCP 端口 & 110（明文）/ 995（SSL） & 143（明文）/ 993（SSL） \\
存储模型 & 下载到本地；是否删除服务器副本由客户端通过协议操作决定 & 以服务器端邮箱状态为准，便于多设备同步 \\
文件夹 & 协议模型较简单，不提供 IMAP 式的服务器端多文件夹管理 & 支持服务器端多文件夹、搜索 \\
工作模式 & 下载—删除（或保留） & 联机操作 \\
离线能力 & 下载后可离线阅读 & 客户端可缓存邮件离线阅读，状态同步需要联网 \\
适用场景 & 单设备收邮件 & 多设备同步邮件 \\
\hline
\end{tabularx}
\end{tablebox}

\textbf{POP3 的三个阶段：}
\begin{enumerate}
    \item \textbf{特许阶段（Authorization）}：发送用户名和密码（USER, PASS）。
    \item \textbf{事务阶段（Transaction）}：列出邮件（LIST）、获取邮件（RETR）、删除邮件（DELE）。
    \item \textbf{更新阶段（Update）}：QUIT 命令后服务器执行删除操作。
\end{enumerate}

\textbf{IMAP} 允许用户在服务器上创建文件夹、移动邮件、只下载邮件头等，比 POP3 复杂但更灵活。IMAPS 使用 TLS/SSL 加密。
\end{conclusion}

\subsubsection{MIME — 多用途互联网邮件扩展}

\begin{conclusion}{MIME 扩展}{mime}
由于 SMTP 仅支持 7 位 ASCII 文本，MIME（Multipurpose Internet Mail Extensions）通过在邮件首部增加字段来支持：
\begin{itemize}
    \item 非 ASCII 文本（如中文、日文等国际字符）；
    \item 二进制附件（图片、音频、视频、PDF 等）；
    \item 图文混合的邮件正文（multipart）。
\end{itemize}

\textbf{MIME 新增的首部字段：}
\begin{itemize}
    \item \texttt{MIME-Version: 1.0} — 声明使用 MIME。
    \item \texttt{Content-Type: text/plain; charset="utf-8"} — 内容的 MIME 类型和字符编码。
    \item \texttt{Content-Transfer-Encoding: base64} — 内容传输编码方式（Base64, Quoted-Printable）。
\end{itemize}

MIME 对 SMTP/POP3 协议本身不做改动，仅仅在邮件数据的首部增加了上述字段。发送方用 Base64 等将二进制编码为 7 位 ASCII，接收方解码还原。

\textbf{常见 Content-Type：}
\texttt{text/plain}, \texttt{text/html}, \texttt{image/jpeg}, \texttt{application/pdf}, \texttt{multipart/mixed}（混合多种内容），\texttt{multipart/alternative}（同一内容多种格式）。
\end{conclusion}

\subsection{万维网与 HTTP}

\subsubsection{URL 的基本概念}

\begin{mydef}{统一资源定位符 URL}{url-def}
URL（Uniform Resource Locator）用于唯一定位互联网上的资源：
\[
\texttt{<protocol>://<host>:<port>/<path>}
\]
\textbf{示例：} \texttt{http://www.example.com:80/index.html}
\begin{itemize}
    \item \textbf{协议（protocol）}：http, https, ftp 等。
    \item \textbf{主机（host）}：域名或 IP 地址。
    \item \textbf{端口（port）}：可选，默认 HTTP 为 80，HTTPS 为 443。
    \item \textbf{路径（path）}：服务器上的资源路径。
\end{itemize}
\end{mydef}

\subsubsection{HTTP 报文格式}

\begin{conclusion}{HTTP 请求报文格式}{http-request-format}
HTTP 报文分为\textbf{请求报文}和\textbf{响应报文}，两者结构类似：起始行 + 首部行 + 空行 + 实体主体。

\textbf{请求行（Request Line）：}
\[
\underbrace{\texttt{GET}}_{\text{方法}}\quad\underbrace{\texttt{/index.html}}_{\text{URL}}\quad\underbrace{\texttt{HTTP/1.1}}_{\text{版本}}
\]

\textbf{常见 HTTP 请求方法：}
\begin{tablebox}{HTTP 请求方法}
\centering
\begin{tabular}{|l|l|}
\hline
\textbf{方法} & \textbf{含义} \\
\hline
GET & 请求获取资源（幂等） \\
POST & 向服务器提交数据（如表单） \\
HEAD & 类似 GET 但只返回首部，不返回主体 \\
PUT & 上传/替换指定资源 \\
DELETE & 删除指定资源 \\
OPTIONS & 查询服务器支持的方法 \\
\hline
\end{tabular}
\end{tablebox}

\textbf{首部行（Header Lines）：} 格式为「字段名: 值」，常见：
\begin{itemize}
    \item \texttt{Host: www.example.com} — 目标主机（HTTP/1.1 必须，支持虚拟主机）。
    \item \texttt{Connection: close} — 请求本次响应后关闭连接。HTTP/1.1 默认持久连接；\texttt{keep-alive} 在 HTTP/1.0 中常用于协商持久连接，在 HTTP/1.1 中通常无需靠它启用默认持久性。
    \item \texttt{User-Agent: Mozilla/5.0} — 客户端类型。
    \item \texttt{Accept: text/html} — 客户端接受的内容类型。
    \item \texttt{Content-Length: 348} — 实体主体的字节数。
\end{itemize}
\end{conclusion}

\begin{conclusion}{HTTP 响应报文格式}{http-response-format}
\textbf{状态行（Status Line）：}
\[
\underbrace{\texttt{HTTP/1.1}}_{\text{版本}}\quad\underbrace{\texttt{200 OK}}_{\text{状态码 + 短语}}\quad
\]

\textbf{常见 HTTP 状态码：}
\begin{tablebox}{HTTP 状态码分类}
\centering
\begin{tabular}{|c|l|l|}
\hline
\textbf{状态码} & \textbf{短语} & \textbf{含义} \\
\hline
200 & OK & 请求成功，响应包含请求的资源 \\
301 & Moved Permanently & 永久重定向（URL 已更改） \\
302 & Found & 临时重定向（原来 URL 仍可用） \\
304 & Not Modified & 缓存有效，不需要重新传输 \\
400 & Bad Request & 请求报文有误 \\
401 & Unauthorized & 需要认证 \\
403 & Forbidden & 服务器拒绝请求 \\
404 & Not Found & 请求的资源不存在 \\
500 & Internal Server Error & 服务器内部错误 \\
502 & Bad Gateway & 网关/代理从上游收到无效响应 \\
503 & Service Unavailable & 服务器暂时不可用（过载/维护）\\
\hline
\end{tabular}
\end{tablebox}

\textbf{记忆规律：}
\begin{itemize}
    \item \textbf{2xx}：成功（200 正常，204 无内容，206 部分内容）。
    \item \textbf{3xx}：重定向（301 永久，302 临时，304 缓存有效）。
    \item \textbf{4xx}：客户端错误（400 请求有问题，404 找不到，403 无权限）。
    \item \textbf{5xx}：服务器错误（500 内部错误，503 服务不可用）。
\end{itemize}
\end{conclusion}

\subsubsection{非持久连接 vs 持久连接}

\begin{conclusion}{HTTP 连接的两种方式}{http-connection}
这是\textbf{408 计算题的核心考点}。

\textbf{非持久连接（Non-persistent / HTTP/1.0 默认）：}
每个 TCP 连接只传送一个对象（一个 HTML 文件 + 每个嵌入图片各需一个连接）。过程：
\begin{enumerate}
    \item 建立 TCP 连接（1 RTT）。
    \item 发送 HTTP 请求，接收响应（1 RTT）。
    \item 关闭 TCP 连接。
\end{enumerate}
传送 \textbf{1 个 HTML 文件 + $n$ 个内嵌对象}的时间：
\[
T_{\text{非持久}} = \underbrace{(2\text{RTT} + T_{\text{文件}})}_{\text{HTML文件}} \;+\; \underbrace{n \times (2\text{RTT} + T_{\text{对象}})}_{\text{内嵌对象}}
\]
\textbf{总时间为}：$2(n+1)$ RTT + 各对象传输时间。

若使用并行连接（浏览器同时打开多个 TCP 连接，常见为 6-10 个），则时间减少。非持久连接按顺序串行发送是理论模型，实际有可能并行。

\textbf{持久连接（Persistent / HTTP/1.1 默认）：}
一个 TCP 连接上可以传送多个对象。过程：
\begin{enumerate}
    \item 建立 TCP 连接（1 RTT）。
    \item 传送所有对象（HTML + 所有内嵌对象）依次发送和接收。
    \item 超时后才关闭连接。
\end{enumerate}

若先取得基本 HTML 后才能知道 $n$ 个内嵌对象的 URL，且不使用流水线，则每个对象的请求—响应各占 1 RTT（忽略传输时间）：
\[
T_{\text{持久、非流水}} = \bigl[1+1+n\bigr]\text{RTT} + \sum \text{各对象传输时间}.
\]
其中三个部分依次是建立 TCP、取得基本 HTML、依次取得 $n$ 个内嵌对象。

\textbf{流水线方式（Pipelining）：} 取得基本 HTML 并获知对象 URL 后，客户端可连续发出所有内嵌对象请求，不必等待前一个响应。于是
\[
T_{\text{流水线}} = \bigl[1+1+1\bigr]\text{RTT} + \sum \text{各对象传输时间}.
\]
三个 RTT 依次用于建立 TCP、取得基本 HTML、流水线取得全部内嵌对象。若题目另行假定客户端在建连前已经知道所有对象 URL，才可能省去「先取 HTML」这一 RTT。HTTP 时间题必须服从题设给定的连接并行度、URL 可知时刻和是否忽略传输时间，不能脱离模型套固定数值。
\end{conclusion}

\begin{attention}
\textbf{408 高分技巧 — HTTP 时间计算题：}
\begin{itemize}
    \item \textbf{RTT（往返时间）}的定义：一个分组从客户端到服务器再返回的时间。
    \item TCP 三次握手需要 1 RTT（前两次握手）后才开始发送数据。因此\textbf{建立连接本身消耗 1 RTT}。
    \item 非持久连接，每个对象额外浪费 2 RTT（1 连接建立 + 1 请求—响应）。
    \item \textbf{并行持久连接（HTTP/1.1 常用）}：打开 $k$ 个并行连接，每个连接上持久传输。需要综合考虑并行度和持久性。
    \item 常见题型：给定网页包含一个 HTML 和 $n$ 个图片，使用非持久串行 / 非持久并行 / 持久 / 流水线方式，计算总时间。
\end{itemize}
\end{attention}

\subsubsection{Cookie}

\begin{conclusion}{HTTP Cookie 机制}{cookie}
HTTP 本身是无状态协议，Cookie 是用来\textbf{在无状态协议上保持状态}的机制。

\textbf{Cookie 的四个组件：}
\begin{enumerate}
    \item 在 HTTP 响应报文中的 \texttt{Set-Cookie} 首部行；
    \item 在 HTTP 请求报文中的 \texttt{Cookie} 首部行；
    \item 浏览器端存储的 Cookie 文件；
    \item Web 站点后端数据库中的 Cookie 记录。
\end{enumerate}

\textbf{工作流程：}
\begin{enumerate}
    \item 客户端首次访问服务器（无 Cookie）。
    \item 服务器在响应中加上：\texttt{Set-Cookie: 1678}（服务器生成唯一 ID）。
    \item 浏览器将 Cookie 号存储在本地 Cookie 文件中。
    \item 后续每次请求该站点时，浏览器自动携带 \texttt{Cookie: 1678} 首部。
    \item 服务器根据 Cookie 号在数据库中查到用户的状态信息（购物车、登录态等），从而在无状态的 HTTP 上实现了有状态的会话。
\end{enumerate}

\textbf{Cookie 的属性：} \texttt{Expires}（过期时间），\texttt{Domain}（作用域），\texttt{Path}（路径），\texttt{Secure}（仅 HTTPS），\texttt{HttpOnly}（不可被 JavaScript 访问）。
\end{conclusion}

\subsection{动态主机配置协议 DHCP}

\begin{mydef}{DHCP 协议}{dhcp-def}
DHCP（Dynamic Host Configuration Protocol）用于为主机自动配置 IP 地址、子网掩码、默认网关、DNS 服务器等网络参数。运行在 UDP 端口 67（服务器）和 68（客户端）上。DHCP 是即插即用联网的关键协议。
\end{mydef}

\subsubsection{DHCP 工作过程 — 四步握手}

\begin{conclusion}{DHCP 四步过程}{dhcp-four-steps}
\textbf{这是 408 选择题的常考点，四个步骤的顺序和报文类型必须牢记。}

\begin{tablebox}{DHCP 四步过程}
\centering
\small
\begin{tabularx}{\linewidth}{|c|p{0.25\linewidth}|X|}
\hline
\textbf{步骤} & \textbf{报文} & \textbf{说明} \\
\hline
1 & DHCPDISCOVER & 客户端广播「发现」报文（源 IP 0.0.0.0，目标 255.255.255.255）寻找 DHCP 服务器 \\
2 & DHCPOFFER & 服务器单播/广播「提供」报文，包含建议的 IP 地址、子网掩码、租用期 \\
3 & DHCPREQUEST & 客户端广播「请求」报文，选择某个服务器提供的 IP 地址 \\
4 & DHCPACK & 服务器「确认」报文，正式分配 IP 地址 \\
\hline
\end{tabularx}
\end{tablebox}

\textbf{详细过程：}
\begin{enumerate}
    \item \textbf{DHCPDISCOVER（发现阶段）}：客户端在物理子网内广播 UDP 报文（端口 67），源 IP = 0.0.0.0，目的 IP = 255.255.255.255（受限广播）。网络中可能有多个 DHCP 服务器都会响应。

    \item \textbf{DHCPOFFER（提供阶段）}：每个 DHCP 服务器都回应 DHCPOFFER 报文（包含：可分配的 IP 地址 yiaddr、子网掩码、租用期、DHCP 服务器标识等）。客户端可能收到多个 OFFER。

    \item \textbf{DHCPREQUEST（请求阶段）}：客户端选择一个 OFFER，再次向全网广播 DHCPREQUEST 报文，明确声明选择哪个 DHCP 服务器（报文中包含该服务器的标识符）。\textbf{广播的目的：} 让其他 DHCP 服务器知道自己未被选中，释放预留的 IP。

    \item \textbf{DHCPACK（确认阶段）}：被选中的服务器发送 DHCPACK 确认分配。客户端收到后就可以正式使用该 IP 地址。租用期过半时客户端需要续租（发送 DHCPREQUEST 直接到服务器）。
\end{enumerate}

\textbf{记忆口诀：} \textbf{「D-O-R-A」}——Discover $\to$ Offer $\to$ Request $\to$ ACK。
\end{conclusion}

\begin{attention}
\textbf{408 常见陷阱：}
\begin{itemize}
    \item DHCPDISCOVER 和 DHCPREQUEST 都是广播（因为客户端此时还没有有效 IP），DHCPOFFER 和 DHCPACK 可以是广播或单播。
    \item DHCP 基于 UDP，不是 TCP。DHCP 服务器用 67 端口，客户端用 68 端口。
    \item DHCP 分配的 IP 地址是\underline{临时的}（有租用期），不是永久的。
    \item 若 DHCP 服务器与客户端不在同一子网，需通过 \textbf{DHCP 中继代理}（Relay Agent）转发。
    \item DHCP 是应用层协议，但它的配置直接影响网络层（IP 地址）和传输层的工作。
\end{itemize}
\end{attention}

\subsection{408 高频考点速查}

\begin{conclusion}{应用层 — 408常考点}{cn6-keypoints}
\begin{enumerate}
    \item \textbf{DNS 递归 vs 迭代查询}：典型模型中主机请求本地 DNS 递归解析，本地 DNS 再向层次服务器进行迭代查询；是否提供递归服务仍取决于服务器策略和请求标志，不能把“主机到本地”写成无条件必然递归。
    \item \textbf{FTP 控制连接 (21) 和数据连接 (20)}：主动模式（服务器端口 20 发起）vs 被动模式（客户端主动连接服务器随机端口）。
    \item \textbf{SMTP/POP3/IMAP 功能区分}：SMTP 只管「发」，POP3/IMAP 只管「收」。IMAP 多设备同步。
    \item \textbf{HTTP 持久连接的时间计算}：RTT 的概念，非持久每个对象 $+2$RTT，持久一个 TCP 连接内传送所有对象。
    \item \textbf{DHCP 四步过程 D-O-R-A}：DISCOVER(广播) $\to$ OFFER $\to$ REQUEST(广播) $\to$ ACK。
    \item \textbf{HTTP 状态码分类}：$2\text{xx}$成功、$3\text{xx}$重定向、$4\text{xx}$客户端错误、$5\text{xx}$服务器错误。
    \item \textbf{Cookie 机制}：无状态 HTTP 下保持会话状态的 4 个组件。
    \item \textbf{MIME 的作用}：弥补 SMTP 仅支持 7 位 ASCII 的不足。
\end{enumerate}
\end{conclusion}

\newpage
\section{习题}

\subsection{基础过关题}

\begin{enumerate}

\item \textbf{（真题风格·选择题）} FTP客户和服务器间传递FTP命令时，使用的连接是
\begin{center}
\begin{tabular}{ll}
(A) 建立在TCP之上的控制连接 & (B) 建立在TCP之上的数据连接 \\[6pt]
(C) 建立在UDP之上的控制连接 & (D) 建立在UDP之上的数据连接
\end{tabular}
\end{center}

\ifshowsolutions\par\noindent\textbf{解析：} FTP使用两个TCP连接：控制连接（端口21）传递命令，数据连接（端口20）传递文件数据。\boxed{\textbf{答案：}(A)}\fi

\vspace{8pt}

\item \textbf{（真题风格·选择题）} 如果本地域名服务器无缓存，当采用递归方法解析另一网络中的某主机域名时，用户主机和本地域名服务器分别发送的域名请求消息条数分别为
\begin{center}
\begin{tabular}{ll}
(A) 1条,1条 & (B) 1条,多条 \\[6pt]
(C) 多条,1条 & (D) 多条,多条
\end{tabular}
\end{center}

\ifshowsolutions\par\noindent\textbf{解析：} 递归解析：用户主机向本地域名服务器发1条请求，本地域名服务器代替用户完成全部查询（可能发多条请求），最终返回结果。用户发1条，本地服务器发多条。\boxed{\textbf{答案：}(B)}\fi

\vspace{8pt}

\item \textbf{（真题风格·选择题）} 电子邮件依次经历三个阶段：（1）用户1向其邮件服务器提交邮件；（2）发送方邮件服务器把邮件传给接收方邮件服务器；（3）用户2从接收方邮件服务器读取邮件。三个阶段分别使用的协议最可能是
\begin{center}
\begin{tabular}{ll}
(A) SMTP, SMTP, POP3 & (B) SMTP, POP3, SMTP \\[6pt]
(C) POP3, SMTP, SMTP & (D) SMTP, SMTP, SMTP
\end{tabular}
\end{center}

\ifshowsolutions\par\noindent\textbf{解析：} 用户1$\to$服务器1:SMTP；服务器1$\to$服务器2:SMTP（邮件服务器之间也使用SMTP）；服务器2$\to$用户2:POP3（用户从服务器取邮件）。\boxed{\textbf{答案：}(A)}\fi

\vspace{8pt}

\item \textbf{（真题风格·选择题）} 使用浏览器访问某大学Web网站主页时，不可能使用到的协议是
\begin{center}
\begin{tabular}{ll}
(A) PPP & (B) ARP \\[6pt]
(C) SMTP & (D) HTTP
\end{tabular}
\end{center}

\ifshowsolutions\par\noindent\textbf{解析：} 访问Web使用HTTP协议，需要ARP解析MAC地址，可能使用PPP（若是拨号连接）。SMTP是电子邮件发送协议，与Web访问无关。\boxed{\textbf{答案：}(C)}\fi

\vspace{8pt}

\item \textbf{（真题风格·选择题）} 在TCP/IP体系结构中，以下协议中属于应用层的是
\begin{center}
\begin{tabular}{ll}
(A) ARP & (B) ICMP \\[6pt]
(C) DHCP & (D) TCP
\end{tabular}
\end{center}

\ifshowsolutions\par\noindent\textbf{解析：} ARP属网络接口层，ICMP属网际层，TCP属传输层，DHCP属应用层（基于UDP）。\boxed{\textbf{答案：}(C)}\fi

\end{enumerate}

\subsection{真题风格与综合训练}

\begin{enumerate}[resume]

\item \textbf{（真题风格·综合题）} DNS解析过程有递归和迭代两种方式。某主机需要解析域名\texttt{www.example.com}，根域名服务器、\texttt{.com}顶级域名服务器、\texttt{example.com}权威域名服务器依次参与。
\begin{enumerate}
    \item[(1)] 分别画出递归解析和迭代解析的查询消息流图；
    \item[(2)] 哪种方式对本地域名服务器的负担更大？哪种方式对根域名服务器的负担更大？
\end{enumerate}

\ifshowsolutions\par\noindent\textbf{解析：}
(1) 完全递归：主机$\to$本地 DNS$\to$根 DNS$\to$.com DNS$\to$example.com 权威 DNS；结果沿相反方向逐级返回。迭代：主机先递归请求本地 DNS；本地 DNS 分别查询根 DNS、.com DNS 和权威 DNS，前两者返回下一层服务器地址，权威 DNS 返回最终记录，本地 DNS 再答复主机。\\
(2) 完全递归要求上级服务器代请求者继续查询，增加上级（包括根服务器）的状态与处理负担；迭代把逐级查询工作交给本地 DNS，因此本地 DNS 的查询负担更大。实际 Internet 中通常是「主机到本地 DNS 递归、本地 DNS 向层次服务器迭代」。\fi

\vspace{8pt}

\item \textbf{（真题风格·综合题）} 简述 HTTP 持久连接和非持久连接的区别。某网页包含 1 个基本 HTML 文件和 5 个 JPEG 图像，所有文件均在同一个服务器上。假定无 DNS 查询时间，TCP 建连后才发送 HTTP 请求，必须先收到基本 HTML 才能知道 5 个图像的 URL，忽略对象传输时间；非持久连接串行建立，持久连接对 5 个图像采用流水线。分别计算两种方式所需的 RTT 数。

\ifshowsolutions\par\noindent\textbf{解析：}
非持久连接：每个对象需单独建立 TCP 连接。HTML 为 2 RTT（建连 1 RTT、请求—响应 1 RTT）；5 个图像各为 2 RTT。串行总计 $6\times 2=12$ RTT。

持久连接流水线：建立唯一的 TCP 连接用 1 RTT；请求并收到基本 HTML 用 1 RTT；得知图像 URL 后一次流水线发出 5 个请求并收到响应，用 1 RTT。总计 $1+1+1=3$ RTT。\fi

\end{enumerate}

\vspace{8pt}
\noindent\textbf{拓展阅读：}
\begin{itemize}
    \item Kurose \& Ross 第2章——DNS、HTTP、FTP、SMTP/POP3的完整协议细节和抓包示例。
    \item 谢希仁 第6章——国内标准术语。
\end{itemize}
